% c11 原书 p.133–148（PDF p.140–155）
设 $(X,\cO_X)$ 是环层空间。由 \ref{it:2.1.6}，我们可以谈论 $\cO_X$-模范畴上左正合共变函子的右导出函子。下面是一些例子。

\noindent\textbf{例}
\begin{enumerate}[label=\arabic*.,leftmargin=2.4em,itemsep=0.3em]
  \item 在拓扑空间 $X$ 上阿贝尔群层构成的范畴上，左正合函子 $\cF\mapsto\cF(X)$ 的第 $i$ 个右导出函子记作 $\cF\mapsto H^i(X,\cF)$。$H^i(X,\cF)$ 称为 $\cF$ 的\emph{第 $i$ 个上同调群}\idx{shangtongdiaoqun}{上同调群}。对 $X$ 的每个开子集 $U$，定义 $H^i(U,-)$ 为左正合函子 $\cF\mapsto\cF(U)$ 的第 $i$ 个右导出函子。我们常把 $H^0(U,\cF)$ 记作 $\Gamma(U,\cF)$。我们有 $\Gamma(U,\cF)=\cF(U)$。

  即使 $(X,\cO_X)$ 是环层空间、$\cF$ 是 $\cO_X$-模，我们也把 $\cF$ 看作阿贝尔群层来计算 $H^i(X,\cF)$。此时 $H^i(X,\cF)$ 有 $\cO_X(X)$-模结构。我们将在后面看到，这样定义的函子 $H^i(X,-)$ 也是 $\cO_X$-模范畴上 $\Gamma(X,-)$ 的第 $i$ 个导出函子。
  \item 设 $(X,\cO_X)$ 是环层空间，$\cF$ 是 $\cO_X$-模。在 $\cO_X$-模范畴上，左正合函子 $\cG\mapsto\Hom_{\cO_X}(\cF,\cG)$ 的第 $i$ 个右导出函子记作 $\cG\mapsto\Ext^i_{\cO_X}(\cF,\cG)$；左正合函子 $\cG\mapsto\mathcal{H}om_{\cO_X}(\cF,\cG)$ 的第 $i$ 个右导出函子记作 $\cG\mapsto\mathcal{E}xt^i_{\cO_X}(\cF,\cG)$。
  \item 设 $f\colon X\to Y$ 是拓扑空间之间的连续映射。在阿贝尔群层范畴上，左正合函子 $\cF\mapsto f_*\cF$ 的第 $i$ 个导出函子记作 $\cF\mapsto\R^if_*\cF$，称为 $\cF$ 的\emph{第 $i$ 个高阶直接像}\idx{gaojiezhijiexiang}{高阶直接像}。即使 $(f,f^\#)\colon(X,\cO_X)\to(Y,\cO_Y)$ 是环层空间的态射、$\cF$ 是 $\cO_X$-模，我们也把 $\cF$ 看作阿贝尔群层来计算 $\R^if_*\cF$。此时 $\R^if_*\cF$ 有 $\cO_Y$-模结构。我们将在后面看到，这样定义的函子 $\R^if_*$ 也是 $\cO_X$-模范畴上 $f_*$ 的第 $i$ 个右导出函子。
  \item 设 $A$ 是环。对任意 $A$-模 $M$，我们都能找到满态射 $F\to M$ 使 $F$ 是自由 $A$-模。特别地，$A$-模范畴有足够投射对象。所以我们可以谈论 $A$-模范畴上右正合共变函子的左导出函子。对任意 $A$-模 $M$，右正合函子 $N\mapsto M\tensor N$ 的第 $i$ 个左导出函子记作 $N\mapsto\Tor_i^A(M,N)$。给定 $N$ 的任意投射分解
  \[
  \cdots\to P^{-1}\to P^0\to N\to 0,
  \]
  我们有
  \[
  \Tor_i^A(M,N)\cong H^{-i}(M\tensor_A P^\cdot).
  \]
\end{enumerate}
\begin{tnote}
这里 $\Tor_i^A(M,N)\cong H^{-i}(M\tensor_AP^\cdot)$ 用的是上同调而非同调，是因为书中的投射分解按上链编号（$\cdots\to P^{-1}\to P^0\to0$，见原书 p.119）。若改用通常的同调编号，就该写成 $\Tor_i^A(M,N)\cong H_i(M\tensor_AP_\cdot)$。两者只是指标约定不同。[依据：原书 p.119 分解的编号与 p.134 的公式]
\end{tnote}

\begin{lemma}\label{it:2.1.7}
设 $(X,\cO_X)$ 是环层空间，$\cI$ 是内射 $\cO_X$-模。则对 $X$ 的每个开子集 $U$，$\cI|_U$ 是内射 $\cO_U$-模。
\end{lemma}
\begin{proof}
设 $j\colon U\to X$ 是包含映射。对任意 $\cO_U$-模 $\cG$，定义 $j_!\cG$ 为 $X$ 上由预层
\[
V\mapsto
\begin{cases}
\cG(V), & \text{若 } V\subset U,\\
0, & \text{其他情形}
\end{cases}
\]
相伴的层。则 $\cG\mapsto j_!\cG$ 是正合函子。$j_!$ 与 $j^*$ 左伴随，其证明留给读者。对任意 $\cO_U$-模 $\cG$，我们有
\[
\Hom_{\cO_U}(\cG,\cI|_U)\cong\Hom_{\cO_X}(j_!\cG,\cI).
\]
由于 $j_!$ 是正合函子且 $\cI$ 是内射对象，$\Hom_{\cO_X}(j_!-,\cI)$ 是正合函子。所以 $\Hom_{\cO_U}(-,\cI|_U)$ 是正合函子，从而 $\cI|_U$ 是内射 $\cO_U$-模。
\end{proof}

\begin{corollary}\label{it:2.1.8}
设 $X$ 是拓扑空间，$\cF$ 是 $X$ 上阿贝尔群层。则对 $X$ 的每个开子集 $U$ 与每个 $i$，我们有 $H^i(U,\cF)\cong H^i(U,\cF|_U)$。
\end{corollary}
\begin{proof}
设 $\cI^\cdot$ 是 $\cF$ 的内射分解。则 $\cI^\cdot|_U$ 是 $\cF|_U$ 的内射分解。所以我们有
\[
H^i(U,\cF|_U)\cong H^i(\cI^\cdot|_U(U))\cong H^i(\cI^\cdot(U))\cong H^i(U,\cF).
\]
\end{proof}

\begin{proposition}\label{it:2.1.9}
设 $f\colon X\to Y$ 是连续映射。则对 $X$ 上每个阿贝尔群层 $\cF$ 与每个 $i$，$\R^if_*\cF$ 是由预层 $V\mapsto H^i(f^{-1}(V),\cF)$ 相伴的层。
\end{proposition}
\begin{proof}
设 $\cI^\cdot$ 是 $\cF$ 的内射分解。则
\[
\R^if_*\cF=\Ker(f_*\cI^i\to f_*\cI^{i+1})/\Image(f_*\cI^{i-1}\to f_*\cI^i).
\]
所以 $\R^if_*\cF$ 是由预层
\[
\begin{aligned}
V\mapsto{}&\Ker\bigl((f_*\cI^i)(V)\to(f_*\cI^{i+1})(V)\bigr)/\Image\bigl((f_*\cI^{i-1})(V)\to(f_*\cI^i)(V)\bigr)\\
={}&\Ker\bigl(\cI^i(f^{-1}(V))\to\cI^{i+1}(f^{-1}(V))\bigr)/\Image\bigl(\cI^{i-1}(f^{-1}(V))\to\cI^i(f^{-1}(V))\bigr)\\
\cong{}&H^i(f^{-1}(V),\cF)
\end{aligned}
\]
相伴的层。
\end{proof}

拓扑空间 $X$ 上的阿贝尔群层 $\cF$ 称为\emph{松弛的}\idx{songchiceng}{松弛层}，若对 $X$ 的每个非空开子集 $U$，限制映射 $\cF(X)\to\cF(U)$ 是满射。这等价于说：对 $X$ 的任意非空开子集 $V\subset U$，限制映射 $\cF(U)\to\cF(V)$ 是满射。对 $X$ 上任意层 $\cF$，定义层 $C^0(\cF)$，使对 $X$ 的每个开子集 $U$，有
\[
C^0(\cF)(U)=\Bigl\{s\colon U\to\prod_{P\in X}\cF_P \Bigm| s(P)\in\cF_P\Bigr\}.
\]

注意我们有
\[
C^0(\cF)\cong\prod_{P\in X}j_{P*}j_P^*\cF,
\]
其中 $j_P\colon\{P\}\to X$ 是任意 $P\in X$ 处的典范嵌入。显然 $C^0(\cF)$ 是松弛层，且有典范单态射
\[
\cF\to C^0(\cF),
\]
它把每个截面 $s\in\cF(U)$ 映为截面
\[
U\to\prod_{P\in X}\cF_P,\qquad P\mapsto s_P
\]
（属于 $C^0(\cF)(U)$）。令 $C^{-1}(\cF)=\cF$。设 $n\ge0$，并假设已对一切 $i\le n$ 定义层 $C^i(\cF)$ 与态射 $C^{i-1}(\cF)\to C^i(\cF)$。定义
\[
C^{n+1}(\cF)=C^0\bigl(\coker(C^{n-1}(\cF)\to C^n(\cF))\bigr),
\]
并把 $C^n(\cF)\to C^{n+1}(\cF)$ 定义为复合
\[
C^n(\cF)\to\coker(C^{n-1}(\cF)\to C^n(\cF))\to C^0\bigl(\coker(C^{n-1}(\cF)\to C^n(\cF))\bigr).
\]
于是我们有正合序列
\[
0\to\cF\to C^0(\cF)\to C^1(\cF)\to\cdots.
\]
所以 $C^\cdot(\cF)$ 是 $\cF$ 的松弛分解。我们称之为 $\cF$ 的 \emph{Godement 分解}。

\begin{proposition}\label{it:2.1.10}
对每个 $i$，$X$ 上阿贝尔群层范畴上的函子 $\cF\mapsto C^i(\cF)$ 是正合的。
\end{proposition}
\begin{proof}
对每个层 $\cF$，把态射 $\cF\to C^0(\cF)$ 记作 $d_{\cF}^{-1}$，把态射 $C^i(\cF)\to C^{i+1}(\cF)$ 记作 $d_{\cF}^i$。我们将对 $i$ 归纳，同时证明函子 $\cF\mapsto C^i(\cF)$ 与 $\cF\mapsto\coker d_{\cF}^{i-1}$ 都是正合的。假设
\[
0\to\cF\to\cG\to\cH\to 0
\]
是层的正合序列。则对任意 $P\in X$，
\[
0\to\cF_P\to\cG_P\to\cH_P\to 0
\]
正合。这说明对 $X$ 的任意开子集 $U$，
\[
0\to C^0(\cF)(U)\to C^0(\cG)(U)\to C^0(\cH)(U)\to 0
\]
正合。因此
\[
0\to C^0(\cF)\to C^0(\cG)\to C^0(\cH)\to 0
\]
正合。这就证明了函子 $\cF\mapsto C^0(\cF)$ 是正合的。我们有交换图
\[\begin{tikzcd}[column sep=2.4em,row sep=1.8em]
0 \arrow[r] & \cF \arrow[r] \arrow[d,"d_{\cF}^{-1}"'] & \cG \arrow[r] \arrow[d,"d_{\cG}^{-1}"'] & \cH \arrow[r] \arrow[d,"d_{\cH}^{-1}"'] & 0 \\
0 \arrow[r] & C^0(\cF) \arrow[r] & C^0(\cG) \arrow[r] & C^0(\cH) \arrow[r] & 0.
\end{tikzcd}\]
注意所有竖直箭头都是单态射。由蛇引理，我们有正合序列
\[
0\to\coker d_{\cF}^{-1}\to\coker d_{\cG}^{-1}\to\coker d_{\cH}^{-1}\to 0.
\]
所以 $\cF\mapsto\coker d_{\cF}^{-1}$ 是正合函子。假设已证明 $\cF\mapsto C^i(\cF)$ 与 $\cF\mapsto\coker d_{\cF}^{i-1}$ 都是正合函子，则 $\cF\mapsto C^{i+1}(\cF)$ 也是正合函子，因为
\[
C^{i+1}(\cF)=C^0\bigl(\coker d_{\cF}^{i-1}\bigr).
\]
对每个短正合序列
\[
0\to\cF\to\cG\to\cH\to 0,
\]
我们有交换图
\[\begin{tikzcd}[column sep=3.4em,row sep=1.8em]
0 \arrow[r] & \coker d_{\cF}^{i-1} \arrow[r] \arrow[d] & \coker d_{\cG}^{i-1} \arrow[r] \arrow[d] & \coker d_{\cH}^{i-1} \arrow[r] \arrow[d] & 0 \\
0 \arrow[r] & C^{i+1}(\cF) \arrow[r] & C^{i+1}(\cG) \arrow[r] & C^{i+1}(\cH) \arrow[r] & 0.
\end{tikzcd}\]
所有竖直箭头都是单态射。由蛇引理，
\[
0\to\coker d_{\cF}^i\to\coker d_{\cG}^i\to\coker d_{\cH}^i\to 0
\]
正合。所以 $\cF\mapsto\coker d_{\cF}^i$ 是正合函子。
\end{proof}
\begin{tnote}
Godement 分解的构造只用到「茎」与「松弛层」，不依赖 $\cF$ 的任何额外选取，因此 $\cF\mapsto C^\cdot(\cF)$ 是函子性的。2.1.10 的证明正是要说明：$\cF\mapsto C^i(\cF)$ 与 $\cF\mapsto\coker d_{\cF}^{i-1}$ 对 $i$ 同时正合。配合 2.1.11 给出的「松弛层对 $\Gamma(X,-)$ 无环」，Godement 分解可以直接用来计算 $H^i(X,\cF)$。[依据：原书 p.136--139]
\end{tnote}

\begin{lemma}\label{it:2.1.11}
设
\[
0\to\cF'\to\cF\to\cF''\to 0
\]
是拓扑空间 $X$ 上阿贝尔群层的正合序列。若 $\cF'$ 是松弛层，则序列
\[
0\to\cF'(X)\to\cF(X)\to\cF''(X)\to 0
\]
正合，且对每个连续映射 $f\colon X\to Y$，
\[
0\to f_*\cF'\to f_*\cF\to f_*\cF''\to 0
\]
也正合。若 $\cF'$ 与 $\cF$ 都是松弛层，则 $\cF''$ 也是松弛层。
\end{lemma}
\begin{proof}
设 $\cF'$ 是松弛层。设 $s''\in\cF''(X)$ 是截面。由于 $\cF\to\cF''$ 是满射，存在 $X$ 的开覆盖 $\{U_i\}_{i\in I}$，使每个截面 $s''|_{U_i}\in\cF''(U_i)$ 都能提升为 $\cF(U_i)$ 中的截面 $s_i$。考虑截面
\[
s_{ij}=s_j|_{U_i\cap U_j}-s_i|_{U_i\cap U_j}.
\]
它们在态射 $\cF\to\cF''$ 下的像为 $0$。但 $\cF'$ 是这个态射的核。所以我们可以把 $s_{ij}$ 看作 $\cF'(U_i\cap U_j)$ 中的截面。我们断言存在截面 $t_i'\in\cF'(U_i)$ 使
\[
s_{ij}=t_j'|_{U_i\cap U_j}-t_i'|_{U_i\cap U_j}.
\]
我们于是有
\[
(s_i-t_i')|_{U_i\cap U_j}=(s_j-t_j')|_{U_i\cap U_j}.
\]
所以我们可以把 $s_i-t_i'$ 粘成 $\cF(X)$ 中的截面。它被映为截面 $s''\in\cF''(X)$。故 $\cF(X)\to\cF''(X)$ 是满射。容易验证
\[
0\to\cF'(X)\to\cF(X)\to\cF''(X)\to 0
\]
正合。由于对 $X$ 的每个开子集 $U$，$\cF'|_U$ 是松弛层，
\[
0\to\cF'(U)\to\cF(U)\to\cF''(U)\to 0
\]
也正合。特别地，对每个连续映射 $f\colon X\to Y$ 与 $Y$ 的每个开子集 $V$，序列
\[
0\to(f_*\cF')(V)\to(f_*\cF)(V)\to(f_*\cF'')(V)\to 0
\]
正合。所以
\[
0\to f_*\cF'\to f_*\cF\to f_*\cF''\to 0
\]
正合。

现在证明我们的断言。令 $\cS$ 为所有截面族 $(t_i')_{i\in J}$ 的集合，其中 $J\subset I$，$t_i'\in\cF'(U_i)$，且对每对 $i,j\in J$ 有 $s_{ij}=t_j'|_{U_i\cap U_j}-t_i'|_{U_i\cap U_j}$。由 Zorn 引理，$\cS$ 有极大元。只需证明 $\cS$ 中每个极大元 $(t_i')_{i\in J}$ 都满足 $J=I$。设 $J\ne I$，取 $i_0\in I-J$。对每对 $i,j\in J$，在 $U_i\cap U_j\cap U_{i_0}$ 上有
\[
t_j'-t_i'=s_{ij}=s_{i_0 j}-s_{i_0 i}.
\]
因此
\[
t_i'-s_{i_0 i}=t_j'-s_{i_0 j}
\]
在 $U_i\cap U_j\cap U_{i_0}$ 上成立。由于 $U_i\cap U_{i_0}$（$i\in J$）构成集合 $(\bigcup_{i\in J}U_i)\cap U_{i_0}$ 的开覆盖，我们可以把截面 $(t_i'-s_{i_0 i})|_{U_i\cap U_{i_0}}$ 粘成 $\cF'((\bigcup_{i\in J}U_i)\cap U_{i_0})$ 中的截面。由于 $\cF'$ 是松弛层，这个截面可延拓为截面 $t_{i_0}'\in\cF'(U_{i_0})$。但这时 $(t_i')_{i\in J\cup\{i_0\}}$ 定义了 $\cS$ 中的一个元素。这与 $(t_i')_{i\in J}$ 的极大性矛盾。故必有 $J=I$。这证明了我们的断言。

最后设 $\cF'$ 与 $\cF$ 都是松弛层。对 $X$ 的每个开子集 $U$，我们有下述交换图：
\[\begin{tikzcd}[column sep=2.6em,row sep=1.8em]
0 \arrow[r] & \cF'(X) \arrow[r] \arrow[d] & \cF(X) \arrow[r] \arrow[d] & \cF''(X) \arrow[r] \arrow[d] & 0 \\
0 \arrow[r] & \cF'(U) \arrow[r] & \cF(U) \arrow[r] & \cF''(U) \arrow[r] & 0.
\end{tikzcd}\]
我们已证明每行都正合。前两个竖直箭头是满射。所以最后一个竖直箭头是满射，$\cF''$ 是松弛层。
\end{proof}

由 \ref{it:2.1.3} 与 \ref{it:2.1.11}，在阿贝尔群层范畴上，松弛层对函子 $\cF\mapsto\cF(X)$ 是无环的。所以由 \ref{it:2.1.2}，我们可以用松弛分解、特别是 $\cF$ 的 Godement 分解来计算 $H^i(X,\cF)$。若 $X$ 是环层空间，则由 \ref{it:2.1.3} 与 \ref{it:2.1.11}，每个内射 $\cO_X$-模都是松弛层。因此 $\cO_X$-模范畴中的每个内射分解也是松弛分解。所以 $\cO_X$-模范畴上函子 $\cF\mapsto\cF(X)$ 的第 $i$ 个导出函子与 $\cF\mapsto H^i(X,\cF)$ 一致。

类似地，对每个连续映射 $f\colon X\to Y$ 与每个阿贝尔群层 $\cF$，我们可以用松弛分解、特别是 $\cF$ 的 Godement 分解来计算 $\R^if_*\cF$。若 $f\colon X\to Y$ 是环层空间的态射，则函子 $\cF\mapsto\R^if_*\cF$ 是 $\cO_X$-模范畴上 $f_*$ 的第 $i$ 个导出函子。

\begin{proposition}\label{it:2.1.12}
设 $X$ 是诺特拓扑空间。
\begin{enumerate}[label=(\roman*),leftmargin=2.4em]
  \item 设 $\cF_i$（$i\in I$）是 $X$ 上的一族阿贝尔群层。则对所有 $n$，有 $H^n(X,\bigoplus_{i\in I}\cF_i)\cong\bigoplus_{i\in I}H^n(X,\cF_i)$。
  \item 设 $I$ 是有向集，$(\cF_i)_{i\in I}$ 是 $X$ 上阿贝尔群层的有向系统。则对所有 $n$，有 $H^n(X,\dlim_i\cF_i)\cong\dlim_i H^n(X,\cF_i)$。
\end{enumerate}
\end{proposition}
\begin{proof}
注意 (i) 由 (ii) 推出。事实上，层的直和是有限多项直和的有向极限，而 (i) 对有限直和成立，因为 $H^n(X,-)$ 是加性函子，从而与有限直和交换。我们来证明 (ii)。由于 $X$ 是诺特拓扑空间，预层 $U\mapsto\dlim_i\cF_i(U)$ 是层，且与 $\dlim_i\cF_i$ 一致。设 $C^\cdot(\cF_i)$ 是 $\cF_i$ 的 Godement 分解。则 $(C^\cdot(\cF_i))_{i\in I}$ 是层的复形的有向系统。可以证明 $\dlim_i C^\cdot(\cF_i)$ 是 $\dlim_i\cF_i$ 的松弛分解。所以我们有
\[
\begin{aligned}
H^n(X,\dlim_i\cF_i)&\cong H^n(\Gamma(X,\dlim_iC^\cdot(\cF_i)))\cong H^n(\dlim_i\Gamma(X,C^\cdot(\cF_i)))\\
&\cong\dlim_iH^n(\Gamma(X,C^\cdot(\cF_i)))\cong\dlim_iH^n(X,\cF_i).
\end{aligned}
\]
\end{proof}

我们把下面命题的证明留给读者。

\begin{proposition}\label{it:2.1.13}
设 $(X,\cO_X)$ 是环层空间，$\cF$ 与 $\cG$ 是 $\cO_X$-模。则对每个 $i$，$\mathcal{E}xt^i_{\cO_X}(\cF,\cG)$ 是由预层 $U\mapsto\Ext^i_{\cO_U}(\cF|_U,\cG|_U)$ 相伴的层。
\end{proposition}

\begin{proposition}\label{it:2.1.14}
设 $(X,\cO_X)$ 是环层空间，$\cF$ 是 $\cO_X$-模，且
\[
0\to\cG'\to\cG\to\cG''\to 0
\]
是 $\cO_X$-模的正合序列。则我们有长正合序列
\[
\cdots\to\Ext^i_{\cO_X}(\cF,\cG')\to\Ext^i_{\cO_X}(\cF,\cG)\to\Ext^i_{\cO_X}(\cF,\cG'')\to\Ext^{i+1}_{\cO_X}(\cF,\cG')\to\cdots
\]
与
\[
\cdots\to\Ext^i_{\cO_X}(\cG'',\cF)\to\Ext^i_{\cO_X}(\cG,\cF)\to\Ext^i_{\cO_X}(\cG',\cF)\to\Ext^{i+1}_{\cO_X}(\cG'',\cF)\to\cdots.
\]
对 $\mathcal{E}xt_{\cO_X}(-,-)$ 也有类似的长正合序列。
\end{proposition}
\begin{proof}
第一个长正合序列来自 $(\Ext^i(\cF,-))$ 是上同调函子这一事实。设 $\cI^\cdot$ 是 $\cF$ 的内射分解。由于对每个内射 $\cO_X$-模 $\cI$，$\Hom_{\cO_X}(-,\cI)$ 是正合函子，我们有复形的短正合序列
\[
0\to\Hom_{\cO_X}(\cG'',\cI^\cdot)\to\Hom_{\cO_X}(\cG,\cI^\cdot)\to\Hom_{\cO_X}(\cG',\cI^\cdot)\to 0.
\]
第二个长正合序列是与这个复形短正合序列相伴的上同调群的长正合序列。
\end{proof}

设 $\cF$ 与 $\cG$ 是环层空间 $(X,\cO_X)$ 上的 $\cO_X$-模。$\cF$ 被 $\cG$ 的\emph{扩张}是指形如
\[
0\to\cG\to\cE\to\cF\to 0
\]
的 $\cO_X$-模短正合序列。两个扩张
\[
0\to\cG\to\cE_1\to\cF\to 0
\]
与
\[
0\to\cG\to\cE_2\to\cF\to 0
\]
称为同构的，若存在交换图
\[\begin{tikzcd}[column sep=2.6em,row sep=1.6em]
0 \arrow[r] & \cG \arrow[r] \arrow[d,equal] & \cE_1 \arrow[r] \arrow[d] & \cF \arrow[r] \arrow[d,equal] & 0 \\
0 \arrow[r] & \cG \arrow[r] & \cE_2 \arrow[r] & \cF \arrow[r] & 0.
\end{tikzcd}\]
注意中间的竖直箭头必是同构。

\begin{proposition}\label{it:2.1.15}
设 $(X,\cO_X)$ 是环层空间，$\cF$ 与 $\cG$ 是两个 $\cO_X$-模。则 $\Ext^1_{\cO_X}(\cF,\cG)$ 与 $\cF$ 被 $\cG$ 的扩张的同构类构成的集合之间有一一对应。
\end{proposition}
\begin{proof}
固定 $\cG$ 的内射分解
\[
0\to\cG\to\cI^0\xrightarrow{d^0}\cI^1\xrightarrow{d^1}\cdots.
\]
我们有
\[
\begin{aligned}
\Ext^1_{\cO_X}(\cF,\cG)
&\cong\Ker\bigl(\Hom_{\cO_X}(\cF,\cI^1)\to\Hom_{\cO_X}(\cF,\cI^2)\bigr)\\
&\qquad\Big/\Image\bigl(\Hom_{\cO_X}(\cF,\cI^0)\to\Hom_{\cO_X}(\cF,\cI^1)\bigr)\\
&\cong\Hom_{\cO_X}(\cF,Z^1(\cI))\big/\Image\bigl(\Hom_{\cO_X}(\cF,\cI^0)\to\Hom_{\cO_X}(\cF,Z^1(\cI))\bigr).
\end{aligned}
\]
设 $\cS$ 为 $\cF$ 被 $\cG$ 的扩张的同构类构成的集合。定义映射
\[
\Phi\colon\Ext^1_{\cO_X}(\cF,\cG)\to\cS
\]
如下：给定元素
\[
e\in\Ext^1_{\cO_X}(\cF,\cG)\cong\Hom_{\cO_X}(\cF,Z^1(\cI))\big/\Image\bigl(\Hom_{\cO_X}(\cF,\cI^0)\to\Hom_{\cO_X}(\cF,Z^1(\cI))\bigr),
\]
设 $e'\colon\cF\to Z^1(\cI)$ 是 $\Hom_{\cO_X}(\cF,Z^1(\cI))$ 中在 $\Ext^1_{\cO_X}(\cF,\cG)$ 里的像为 $e$ 的元素。我们定义 $\Phi(e)$ 为短正合序列
\[
0\to\cG\to\cI^0\to Z^1(\cI)\to 0
\]
沿态射 $e'\colon\cF\to Z^1(\cI)$ 的拉回的同构类；也就是说，$\Phi(e)$ 是扩张
\[
0\to\cG\to\cI^0\oplus_{e'}\cF\to\cF\to 0
\]
的同构类，其中 $\cI^0\oplus_{e'}\cF$ 是 $\cI^0\oplus\cF$ 中由满足 $d^0(s)=e'(t)$ 的截面 $(s,t)$ 构成的子层，态射 $\cG\to\cI^0\oplus_{e'}\cF$ 由 $\cG\to\cI^0$ 与零态射 $0\colon\cG\to\cF$ 定义，态射 $\cI^0\oplus_{e'}\cF\to\cF$ 是投影。我们断言 $\Phi(e)$ 与 $e'\colon\cF\to Z^1(\cI)$ 的选取无关，从而 $\Phi(e)$ 是良定义的。若 $e''\colon\cF\to Z^1(\cI)$ 是另一个在 $\Ext^1_{\cO_X}(\cF,\cG)$ 中的像为 $e$ 的态射，则对某个态射 $f\in\Hom_{\cO_X}(\cF,\cI^0)$ 有 $e'-e''=d^0f$。我们有交换图
\[\begin{tikzcd}[column sep=2.6em,row sep=1.8em]
0 \arrow[r] & \cG \arrow[r] \arrow[d,equal] & \cI^0\oplus_{e'}\cF \arrow[r] \arrow[d] & \cF \arrow[r] \arrow[d,equal] & 0 \\
0 \arrow[r] & \cG \arrow[r] & \cI^0\oplus_{e''}\cF \arrow[r] & \cF \arrow[r] & 0,
\end{tikzcd}\]
其中中间的竖直箭头 $\cI^0\oplus_{e'}\cF\to\cI^0\oplus_{e''}\cF$ 把 $\cI^0\oplus_{e'}\cF$ 的每个截面 $(s,t)$ 映为 $\cI^0\oplus_{e''}\cF$ 的截面 $(s-f(t),t)$。这说明 $e'$ 的拉回与 $e''$ 的拉回同构。因此 $\Phi(e)$ 与 $e'$ 的选取无关。

定义 $\Psi\colon\cS\to\Ext^1_{\cO_X}(\cF,\cG)$ 如下：设 $E\in\cS$ 是扩张
\[
0\to\cG\to\cE\to\cF\to 0
\]
的同构类。由于 $\cI^0$ 是内射对象，我们可以把态射 $\cG\to\cI^0$ 延拓为态射 $\psi\colon\cE\to\cI^0$。设 $\bar\psi\colon\cF\to Z^1(\cI)$ 是使下图交换的唯一态射：
\[\begin{tikzcd}[column sep=2.6em,row sep=1.8em]
0 \arrow[r] & \cG \arrow[r] \arrow[d,equal] & \cE \arrow[r] \arrow[d,"\psi"'] & \cF \arrow[r] \arrow[d,"\bar\psi"] & 0 \\
0 \arrow[r] & \cG \arrow[r] & \cI^0 \arrow[r] & Z^1(\cI) \arrow[r] & 0.
\end{tikzcd}\]
我们把 $\Psi(E)$ 定义为 $\bar\psi\colon\cF\to Z^1(\cI)$ 在
\[
\Ext^1_{\cO_X}(\cF,\cG)\cong\Hom_{\cO_X}(\cF,Z^1(\cI))\big/\Image\bigl(\Hom_{\cO_X}(\cF,\cI^0)\to\Hom_{\cO_X}(\cF,Z^1(\cI))\bigr)
\]
中的像。我们断言 $\Psi(E)$ 与 $\psi\colon\cE\to\cI^0$ 的选取无关。若 $\psi'\colon\cE\to\cI^0$ 是另一个延拓 $\cG\to\cI^0$ 的态射，则 $\psi-\psi'$ 在 $\Image(\cG\to\cE)$ 上为零，从而诱导态射 $\delta\colon\cF\to\cI^0$。我们有 $\bar\psi-\bar\psi'=d^0\delta$。所以 $\bar\psi$ 与 $\bar\psi'$ 在 $\Ext^1_{\cO_X}(\cF,\cG)$ 中的像相同。因此 $\Psi(E)$ 与 $\psi$ 的选取无关。容易验证 $\Psi$ 是良定义的，且 $\Phi$ 与 $\Psi$ 互为逆映射。
\end{proof}
\begin{tnote}
$\Ext^1$ 与扩张的对应是 Yoneda 型对应：$\Phi$ 取拉回，$\Psi$ 取推出（把 $\cG\to\cI^0$ 沿 $\cG\to\cE$ 延拓出去）。证明中两段「与选取无关」的论证，实际上分别验证了：换代表元只差 $d^0f$，于是拉回相差一个同构；换 $\psi$ 只差 $d^0\delta$，于是推出的像在 $\Ext^1$ 中相同。[依据：原书 p.141--142 的构造与两张交换图]
\end{tnote}

我们把下面命题的证明留给读者。

\begin{proposition}\label{it:2.1.16}
设 $A$ 是环，$M$ 是 $A$-模，且
\[
0\to N'\to N\to N''\to 0
\]
是 $A$-模的正合序列。则我们有长正合序列
\[
\cdots\to\Tor^A_{i+1}(M,N'')\to\Tor^A_i(M,N')\to\Tor^A_i(M,N)\to\Tor^A_i(M,N'')\to\cdots
\]
与
\[
\cdots\to\Tor^A_{i+1}(N'',M)\to\Tor^A_i(N',M)\to\Tor^A_i(N,M)\to\Tor^A_i(N'',M)\to\cdots.
\]
\end{proposition}

$A$-模 $M$ 称为\emph{平坦模}\idx{pingtan}{平坦模}，若函子 $M\tensor_A-$ 是正合的。注意每个自由 $A$-模都是平坦的。由于每个投射 $A$-模都是自由 $A$-模的直和项，每个投射 $A$-模也都是平坦的。

\begin{proposition}\label{it:2.1.17}
设 $A$ 是环，$M$ 是 $A$-模。则以下条件等价：
\begin{enumerate}[label=(\roman*),leftmargin=2.4em]
  \item $M$ 是平坦的。
  \item 对每个 $A$-模 $N$ 与每个 $i\ge1$，有 $\Tor^A_i(M,N)=0$。
  \item 对每个 $A$-模 $N$，有 $\Tor^A_1(M,N)=0$。
  \item 对每个有限生成 $A$-模的正合序列
  \[
  0\to N'\to N,
  \]
  序列
  \[
  0\to M\tensor_A N'\to M\tensor_A N
  \]
  正合。
\end{enumerate}
\end{proposition}
\begin{proof}
(i)$\Rightarrow$(ii) 设
\[
\cdots\to P^{-1}\to P^0\to N\to 0
\]
是 $N$ 的投射分解。由于 $M$ 是平坦的，序列
\[
\cdots\to M\tensor_A P^{-1}\to M\tensor_A P^0\to M\tensor_A N\to 0
\]
正合。所以对每个 $i\ge1$，我们有
\[
\Tor_i^A(M,N)\cong H^{-i}(M\tensor_A P^\cdot)=0.
\]
(ii)$\Rightarrow$(iii) 显然。

(iii)$\Rightarrow$(iv) 设 $N/N'$ 是 $N'\to N$ 的余核。我们有短正合序列
\[
0\to N'\to N\to N/N'\to 0,
\]
从而有长正合序列
\[
\cdots\to\Tor_1^A(M,N/N')\to M\tensor_A N'\to M\tensor_A N\to M\tensor_A N/N'\to 0.
\]
由 (iii)，我们有 $\Tor_1^A(M,N/N')=0$。所以上述长正合序列的一部分是
\[
0\to M\tensor_A N'\to M\tensor_A N.
\]
(iv)$\Rightarrow$(i) 由于 $M\tensor_A-$ 是右正合的，只需证明：对 $A$-模的每个正合序列
\[
0\to N'\xrightarrow{\phi} N,
\]
序列
\[
0\to M\tensor_A N'\xrightarrow{\id_M\tensor\phi} M\tensor_A N
\]
正合。设 $\sum_{i=1}^k m_i\tensor n_i\in\Ker(\id_M\tensor\phi)$。则在 $M\times N$ 中，我们可以写
\[
\begin{aligned}
\sum_{i=1}^k(m_i,\phi(n_i))={}&\sum_{i=1}^{k_1}\bigl((x_{1i}+x_{2i},y_i)-(x_{1i},y_i)-(x_{2i},y_i)\bigr)\\
&+\sum_{i=1}^{k_2}\bigl((x_i,y_{1i}+y_{2i})-(x_i,y_{1i})-(x_i,y_{2i})\bigr)\\
&+\sum_{i=1}^{k_3}\bigl((a_iu_i,v_i)-(u_i,a_iv_i)\bigr).
\end{aligned}
\]
设 $N_1$ 是 $N$ 的由 $\phi(n_i)$、$y_i$、$y_{1i}$、$y_{2i}$ 与 $v_i$ 生成的子模，$N_1'$ 是 $N'$ 的由 $n_i$ 生成的子模，$\phi_1\colon N_1'\to N_1$ 是由 $\phi$ 诱导的同态。则
\[
0\to N_1'\xrightarrow{\phi_1}N_1
\]
是有限生成 $A$-模的正合序列，且 $\sum_{i=1}^k m_i\tensor n_i$ 属于 $\Ker(\id_M\tensor\phi_1)$。由 (iv)，$\sum_{i=1}^k m_i\tensor n_i$ 在 $M\tensor_A N_1'$ 中为零，从而在 $M\tensor_A N'$ 中为零。
\end{proof}

由 \ref{it:2.1.2} 与 \ref{it:2.1.17}，我们可以用 $N$ 的平坦分解来计算 $\Tor_i^A(M,N)$。

注意，对环 $A$ 的每个乘性集 $S$，环 $S^{-1}A$ 是平坦 $A$-代数，因为函子
\[
M\mapsto S^{-1}A\tensor_A M\cong S^{-1}M
\]
是 $A$-模范畴上的正合函子。若 $A$ 是诺特的，则对 $A$ 的任何理想 $I$，$A$ 的 $I$-进完全化 $\widehat A$ 是平坦 $A$-代数。事实上，由 \ref{it:1.5.5}，对任意有限生成 $A$-模 $M$ 有
\[
\widehat A\tensor_A M\cong\widehat M,
\]
而由 \ref{it:1.5.4}，函子 $M\mapsto\widehat M$ 在有限生成 $A$-模范畴上是正合的。设 $A\to B$ 是环同态。若 $M$ 是平坦 $A$-模，则 $B\tensor_A M$ 是平坦 $B$-模。

\begin{proposition}\label{it:2.1.18}
设 $A$ 是环，$M$ 是有限表示 $A$-模。则 $M$ 是平坦的当且仅当 $M$ 是投射的。
\end{proposition}
\begin{proof}
我们已经看到，任意投射 $A$-模都是平坦的。设 $M$ 是有限表示的平坦 $A$-模。为证明 $M$ 是投射的，只需证明 $M^\sim$ 是概形 $\Spec A$ 上的局部自由层（由 \ref{it:1.4.7}）。由 \ref{it:1.4.1} (ii)，只需证明对每个 $\fp\in\Spec A$，$M_\fp$ 是自由 $A_\fp$-模。于是我们化归为证明：对每个局部环 $A$，每个有限表示的平坦 $A$-模都是自由的。设 $\fm$ 是 $A$ 的极大理想。选取 $x_1,\dots,x_n\in M$，使它们在 $M/\fm M$ 中的像构成 $A/\fm$-向量空间 $M/\fm M$ 的一组基。考虑同态
\[
A^n\to M,\qquad e_i\mapsto x_i\quad(i=1,\dots,n),
\]
其中对每个 $i$，$e_i$ 是 $A^n$ 中第 $i$ 个分量为 $1$、其余分量为 $0$ 的元素。由 Nakayama 引理，$A^n\to M$ 是满射。设 $R$ 是它的核。我们有正合序列
\[
0\to R\to A^n\to M\to 0.
\]
由于 $M$ 是平坦的，我们有 $\Tor_1^A(M,A/\fm)=0$。与上述短正合序列相伴的 $\Tor_i^A(-,A/\fm)$ 的第二个长正合序列（见 \ref{it:2.1.16}）给出正合序列
\[
0\to R\tensor_A A/\fm\to A^n\tensor_A A/\fm\to M\tensor_A A/\fm\to 0,
\]
即
\[
0\to R/\fm R\to(A/\fm)^n\to M/\fm M\to 0.
\]
按我们对 $A^n\to M$ 的定义，同态 $(A/\fm)^n\to M/\fm M$ 是同构。所以 $R/\fm R=0$。由 $M$ 是有限表示这一条件推出 $R$ 是有限生成的，这一步留给读者。由 Nakayama 引理，我们有 $R=0$。于是 $A^n\to M$ 是同构，从而 $M$ 是自由的。
\end{proof}

\begin{proposition}\label{it:2.1.19}
设 $\phi\colon A\to B$ 是环同态，$M$ 是 $B$-模。则以下条件等价：
\begin{enumerate}[label=(\roman*),leftmargin=2.4em]
  \item $M$ 作为 $A$-模是平坦的。
  \item 对每个 $\fq\in\Spec B$，$M_\fq$ 作为 $A$-模是平坦的。
  \item 对每个 $\fq\in\Spec B$，$M_\fq$ 作为 $A_{\phi^{-1}(\fq)}$-模是平坦的。
\end{enumerate}
\end{proposition}
\begin{proof}
(i)$\Rightarrow$(ii) 对每个 $A$-模 $N$，我们有
\[
M_\fq\tensor_A N\cong B_\fq\tensor_B(M\tensor_A N).
\]
由于 $M$ 作为 $A$-模是平坦的，$A$-模范畴上的函子 $M\tensor_A-$ 是正合的。而 $B_\fq\tensor_B-$ 是 $B$-模范畴上的正合函子。所以 $A$-模范畴上的函子 $M_\fq\tensor_A-$ 是正合的，从而 $M_\fq$ 作为 $A$-模是平坦的。

(ii)$\Rightarrow$(iii) 由于 $M_\fq$ 作为 $A$-模是平坦的，$(M_\fq)_{\phi^{-1}(\fq)}$ 作为 $A_{\phi^{-1}(\fq)}$-模是平坦的。我们有 $(M_\fq)_{\phi^{-1}(\fq)}\cong M_\fq$。

(iii)$\Rightarrow$(i) 我们需要证明 $M\tensor_A-$ 是正合函子。只需证明对每个 $\fq\in\Spec B$，$B_\fq\tensor_BM\tensor_A-$ 是正合函子。对每个 $A$-模 $N$，我们有
\[
B_\fq\tensor_BM\tensor_AN\cong M_\fq\tensor_AN\cong M_\fq\tensor_{A_{\phi^{-1}(\fq)}}A_{\phi^{-1}(\fq)}\tensor_AN\cong M_\fq\tensor_{A_{\phi^{-1}(\fq)}}N_{\phi^{-1}(\fq)}.
\]
我们的断言由「$M_\fq$ 作为 $A_{\phi^{-1}(\fq)}$-模是平坦的」这一假设推出。
\end{proof}

\begin{proposition}\label{it:2.1.20}
设 $A$ 是环，$M$ 是 $A$-模。则以下条件等价：
\begin{enumerate}[label=(\roman*),leftmargin=2.4em]
  \item 一个 $A$-模序列
  \[
  N'\xrightarrow{\phi}N\xrightarrow{\psi}N''
  \]
  正合，当且仅当
  \[
  M\tensor_AN'\xrightarrow{\id_M\tensor\phi}M\tensor_AN\xrightarrow{\id_M\tensor\psi}M\tensor_AN''
  \]
  正合。
  \item $M$ 是平坦 $A$-模，且对每个 $A$-模 $N$，条件 $M\tensor_AN=0$ 蕴含 $N=0$。
  \item $M$ 是平坦 $A$-模，且对每个 $A$-模同态 $\phi\colon N\to N'$，条件 $\id_M\tensor\phi=0$ 蕴含 $\phi=0$。也就是说，函子 $M\tensor_A-$ 是正合且忠实的。
  \item $M$ 是平坦 $A$-模，且对 $A$ 的每个极大理想 $\fm$，有 $\fm M\ne M$。
\end{enumerate}
\end{proposition}

当 $M$ 满足上述条件之一时，我们称 $M$ 是\emph{忠实平坦的}\idx{zhongshipingtan}{忠实平坦} $A$-模。
\begin{proof}
(i)$\Rightarrow$(ii) 假设 $M\tensor_AN=0$。则序列 $0\to M\tensor_AN\to0$ 正合。所以 $0\to N\to0$ 正合。故 $N=0$。

(ii)$\Rightarrow$(iii) 由于 $M$ 是平坦 $A$-模，我们有 $\Image(\id_M\tensor\phi)\cong M\tensor_A\Image\phi$。若 $\id_M\tensor\phi=0$，则 $M\tensor_A\Image\phi=0$。故 $\Image\phi=0$。

(iii)$\Rightarrow$(iv) 假设对 $A$ 的某个极大理想 $\fm$ 有 $\fm M=M$。则 $M\tensor_AA/\fm\cong M/\fm M=0$。所以 $\id_M\tensor\id_{A/\fm}=0$。这说明 $\id_{A/\fm}=0$，从而 $A/\fm=0$。这不可能。

(iv)$\Rightarrow$(ii) 假设 $N\ne0$。设 $x$ 是 $N$ 中的非零元素。考虑同态
\[
A\to N,\qquad a\mapsto ax.
\]
设 $\fa$ 是它的核。我们有 $\fa\ne A$，从而 $\fa$ 含于某个极大理想 $\fm$ 中。由于 $\fm M\ne M$，我们有 $\fa M\ne M$，从而 $M\tensor_AA/\fa\ne0$。由于 $M$ 是平坦的，单态射 $A/\fa\to N$ 诱导单态射 $M\tensor_AA/\fa\to M\tensor_AN$。所以我们有 $M\tensor_AN\ne0$。

(iii)$\Rightarrow$(i) 假设
\[
M\tensor_AN'\xrightarrow{\id_M\tensor\phi}M\tensor_AN\xrightarrow{\id_M\tensor\psi}M\tensor_AN''
\]
正合。则 $\id_M\tensor\psi\phi=0$。所以 $\psi\phi=0$。由于 $M$ 是平坦的，我们有
\[
M\tensor_A(\Ker\psi/\Image\phi)\cong\Ker(\id_M\tensor\psi)/\Image(\id_M\tensor\phi)=0.
\]
故 $\Ker\psi/\Image\phi=0$。
\end{proof}

\begin{proposition}\label{it:2.1.21}
设 $\phi\colon A\to B$ 是环同态。则以下条件等价：
\begin{enumerate}[label=(\roman*),leftmargin=2.4em]
  \item $B$ 是忠实平坦 $A$-代数。
  \item $B$ 是平坦 $A$-代数，且对任意 $A$-模 $M$，同态
  \[
  M\to B\tensor_AM,\qquad x\mapsto1\tensor x
  \]
  是单射。（特别地，$A\to B$ 是单射。）
  \item $B$ 是平坦 $A$-代数，且对 $A$ 的每个理想 $\fa$，有 $\phi^{-1}(\fa B)=\fa$。
  \item $B$ 是平坦 $A$-代数，且对 $A$ 的每个极大理想 $\fm$，存在 $B$ 的极大理想 $\fn$ 使 $\phi^{-1}(\fn)=\fm$。
  \item $B$ 是平坦 $A$-代数，且对 $A$ 的每个素理想 $\fp$，存在 $B$ 的素理想 $\fq$ 使 $\phi^{-1}(\fq)=\fp$。
\end{enumerate}
\end{proposition}
\begin{proof}
(i)$\Rightarrow$(ii) 只需证明 $B\tensor_AM\to B\tensor_A(B\tensor_AM)$ 是单射。事实上它有右逆
\[
B\tensor_A(B\tensor_AM)\to B\tensor_AM,\qquad b_1\tensor(b_2\tensor x)\mapsto b_1b_2\tensor x.
\]
(ii)$\Rightarrow$(iii) 把 (ii) 用于 $M=A/\fa$，我们看到 $\phi$ 诱导单态射 $A/\fa\to B/\fa B$。这说明 $\fa=\phi^{-1}(\fa B)$。
(iii)$\Rightarrow$(iv) 注意 $\fm B\ne B$。否则 $\phi^{-1}(\fm B)=A\ne\fm$。所以存在 $B$ 的极大理想 $\fn$ 使 $\fm B\subset\fn$。我们于是有 $\fm\subset\phi^{-1}(\fn)$。由于 $\fm$ 是极大理想，必有 $\fm=\phi^{-1}(\fn)$。
(iv)$\Rightarrow$(i) 由 \ref{it:2.1.20} (iv) 推出。
(i)$\Rightarrow$(v) 可以证明 $B_\fp$ 是忠实平坦 $A_\fp$-代数。由已证明的 (i) 与 (iv) 的等价性，存在 $B_\fp$ 的极大理想 $\fn$ 使 $\phi_\fp^{-1}(\fn)=\fp A_\fp$，其中 $\phi_\fp\colon A_\fp\to B_\fp$ 是由 $\phi$ 诱导的同态。设 $\fq$ 是 $\fn$ 在典范同态 $B\to B_\fp$ 下的原像。则 $\fq$ 是 $B$ 的素理想且 $\phi^{-1}(\fq)=\fp$。
\end{proof}

\begin{corollary}\label{it:2.1.22}
设 $A\to B$ 是局部环之间的局部同态。若 $B$ 是平坦 $A$-代数，则它是忠实平坦的。
\end{corollary}
\begin{proof}
这由 \ref{it:2.1.21} (iv) 推出。
\end{proof}

\begin{corollary}\label{it:2.1.23}
设 $A$ 是诺特环，$I$ 是 $A$ 的理想，$\widehat A$ 是 $A$ 的 $I$-进完全化。则 $\widehat A$ 是忠实平坦 $A$-代数，当且仅当 $I\subset\operatorname{rad}(A)$。
\end{corollary}
\begin{proof}
我们已经看到 $\widehat A$ 是平坦 $A$-代数。设 $\phi\colon A\to\widehat A$ 是典范同态。假设 $\widehat A$ 是忠实平坦的。则对 $A$ 的每个极大理想 $\fm$，由 \ref{it:2.1.21} (iv)，存在 $\widehat A$ 的极大理想 $\fn$ 使 $\fm=\phi^{-1}(\fn)$。由 \ref{it:1.5.7}，我们有 $\widehat I\subset\fn$。所以
\[
I\subset\phi^{-1}(\widehat I)\subset\phi^{-1}(\fn)=\fm.
\]
于是 $I$ 含于 $A$ 的每个极大理想中。因此 $I\subset\operatorname{rad}(A)$。

反之，假设 $I\subset\operatorname{rad}(A)$。设 $\fm$ 是 $A$ 的极大理想。则 $\fm/I$ 是 $A/I$ 的极大理想。存在 $\widehat A$ 的极大理想 $\fn$，使 $\fm/I$ 通过同构 $A/I\cong\widehat A/\widehat I$ 与 $\fn/\widehat I$ 等同。我们于是有 $\fm=\phi^{-1}(\fn)$。所以由 \ref{it:2.1.21} (iv)，$\widehat A$ 是忠实平坦的。
\end{proof}
\begin{tnote}
2.1.21 的 (v) 与 2.1.20 的 (iv) 一起，刻画了忠实平坦的几何含义：$\Spec B\to\Spec A$ 的像落满全空间，且环上的模在张量之后不会「莫名消失」。2.1.23 则说明：$I$-进完全化只有在 $I\subset\operatorname{rad}(A)$ 时才保持这种性质，这也解释了为何局部环的完全化总是忠实平坦的。[依据：原书 p.146--147，命题 2.1.21 (iv)(v) 与推论 2.1.23]
\end{tnote}

\begin{exercises}
\item \label{ex:2.1.1} 设 $\phi\colon K^\cdot\to L^\cdot$ 是复形态射，如下定义一个复形 $C^\cdot(\phi)$：对每个 $i$，令
\[
C^i(\phi)=L^i\oplus K^{i+1},
\]
并令 $d^i\colon C^i(\phi)\to C^{i+1}(\phi)$ 为态射
\[
L^i\oplus K^{i+1}\to L^{i+1}\oplus K^{i+2},\qquad (x_i,y_{i+1})\mapsto(dx_i+\phi(y_{i+1}),-d(y_{i+1})).
\]
证明 $dd=0$，且我们有上同调对象的长正合序列
\[
\cdots\to H^i(K^\cdot)\to H^i(L^\cdot)\to H^i(C^\cdot(\phi))\to H^{i+1}(K^\cdot)\to\cdots.
\]
由此可知，$C^\cdot(\phi)$ 无环，当且仅当 $\phi$ 对每个 $i$ 都诱导同构 $H^i(K^\cdot)\xrightarrow{\cong}H^i(L^\cdot)$。我们称 $C^\cdot(\phi)$ 为态射 $\phi\colon K^\cdot\to L^\cdot$ 的\emph{映射锥}\idx{yingshezhui}{映射锥}。证明恒等态射 $\id_{K^\cdot}$ 的映射锥 $C^\cdot(\id_{K^\cdot})$ 与零复形同伦。

\item \label{ex:2.1.2} 保持 1 中的假设，并进一步假设我们有复形的短正合序列
\[
0\to K^\cdot\xrightarrow{\phi}L^\cdot\xrightarrow{\psi}M^\cdot\to0.
\]
对每个 $i$，定义 $f\colon C^i(\phi)\to M^i$ 为
\[
f(x_i,y_{i+1})=\psi(x_i).
\]
证明 $f\colon C^\cdot(\phi)\to M^\cdot$ 是复形态射，且它对所有 $i$ 诱导同构
\[
H^i(C^\cdot(\phi))\cong H^i(M^\cdot).
\]
我们有复形的短正合序列
\[
0\to C^\cdot(\id_{K^\cdot})\to C^\cdot(\phi)\xrightarrow{f}M^\cdot\to0.
\]

\item \label{ex:2.1.3} 设 $X$ 是拓扑空间，$Z$ 是 $X$ 的闭子集。对 $X$ 上任意阿贝尔群层 $\cF$，定义 $\Gamma_Z(X,\cF)$ 为 $\cF$ 在 $X$ 上支集含于 $Z$ 的截面构成的群，即那些在 $X-Z$ 上为零的截面。则函子
\[
\cF\mapsto\Gamma_Z(X,\cF)
\]
是左正合的。它的右导出函子记作 $\cF\mapsto H^i_Z(X,\cF)$，称为\emph{支集在 $Z$ 中的上同调群}。设 $Y$ 是 $X$ 的闭子集，$U=X-Y$。则我们有长正合序列
\[
\cdots\to H^i_{Z\cap Y}(X,\cF)\to H^i_Z(X,\cF)\to H^i_{Z\cap U}(U,\cF|_U)\to H^{i+1}_{Z\cap Y}(X,\cF)\to\cdots.
\]

\item \label{ex:2.1.4} 设 $X$ 是拓扑空间，$Z$ 是 $X$ 的闭子集，$\cF$ 是 $X$ 上阿贝尔群层。假设 $\cF$ 是松弛层。则对所有 $i\ge1$，有 $H^i_Z(X,\cF)=0$。所以我们可以用 $\cF$ 的松弛分解来计算 $H^i_Z(X,\cF)$。
\end{exercises}
% ===END===
